% !PRINTED 187
\weberpage{187}
%<<CONT
eines beliebigen, im Raume festen Punktes, bezogen auf dies
Axensystem in der Einheitsstellung $E$, so erh\"alt man die Coordinaten desselben Punktes, bezogen auf das Axensystem in der
Stellung $A$ oder $B$ durch zwei orthogonale Substitutionen $A(x)$
und $B(x)$.

Demnach sind $AB(x)$ die auf das System $A$ bezogenen
Coordinaten eines Punktes $y$ mit den Coordinaten $(y) = B(x)$
im Systeme $E$. Der Punkt $y$ hat also im Systeme $E$ dieselben
Coordinaten, wie der Punkt $x$ im Systeme $B$, d. h. $y$ liegt zur
Einheitsstellung so wie $x$ zur Stellung $B$. F\"uhren wir nun die
Drehung, die von $E$ zu $B$ f\"uhrt, so aus, dass wir den Punkt $x$
festhalten, aber den Punkt $y$ und das Axensystem $A$ die Drehung
begleiten lassen, so gelangt der Punkt $y$ nach $x$, und die Coordinaten von $x$, bezogen auf die neue Lage des Systemes $A$, sind
dieselben, wie die des Punktes $y$ in Bezug auf die urspr\"ungliche
Lage von $A$, d. h. $AB(x)$. Diese neue Lage des Systemes $A$
kann man aber offenbar auch so erreichen, dass man die Drehung,
die zu der Stellung $A$ f\"uhrt, nicht von der Einheitsstellung,
sondern von der Stellung $B$ aus vollzieht.

Aus der Gruppe der eigentlichen erh\"alt man die Gesammtheit
aller orthogonalen Substitutionen durch Zusammensetzung mit
\gesperrt{einer} uneigentlichen, etwa mit $(x_1, x_2, x_3) = (y_1, y_2, - y_3)$,
die, nach der zweiten geometrischen Auffassung, eine Spiegelung
an der Ebene $x_1, x_2$ ist, bei der der Punkt $y$ das Spiegelbild
des Punktes $x$ ist.

Von besonderem Interesse sind nun die in der Gesammtheit
der eigentlichen orthogonalen Substitutionen enthaltenen \gesperrt{endlichen Gruppen}, auf die wir sp\"ater noch n\"aher eingehen werden.
Wir wollen hier nur noch \"uber die geometrische Seite dieser
Frage Folgendes bemerken. Einer solchen endlichen Gruppe $G$
vom Grade $g$ von orthogonalen Substitutionen entspricht eine
endliche Gruppe von Stellungen eines K\"orpers. Denkt man sich
den K\"orper in diesen verschiedenen Stellungen gleichzeitig fixirt
und das Ganze zu einem neuen starren K\"orper vereinigt, so
erh\"alt man ein Gebilde, das in einer endlichen Anzahl $g$ verschiedener Stellungen sich selbst decken kann. Man kann f\"ur
jede Gruppe K\"orper von unendlich vielen verschiedenen Gestalten
finden. Die anschaulichsten und bekanntesten Verh\"altnisse ergeben sich, wenn man den K\"orper von ebenen Fl\"achen begrenzt
annimmt.
